%% =====================================================================
%% Supplementary material body (included by supplement.tex). Every row is generated from the workbook by
%% scripts/gen_from_workbook.py (generated/*.tex); do not edit by hand.
%% =====================================================================
\section{Study-Level Coding of the Included Studies}
\label{app:coding}
Task: Det = detection, Attr = attribution, Loc = localization, Bench = benchmark, Rob = robustness or attack study. Evidence (primary family, Axis B): SPA spatial, FRQ frequency, NOI noise/residual, SRC source/fingerprint, SEM semantic, REC reconstruction, GEO geometry/physical, HYB hybrid. Generators: G = GAN only, D = diffusion only, G+D = both, Multi = multi-family (more than the GAN and diffusion pair). Generalization tests (Axis D): U unseen generator, F cross-family, Ds cross-dataset, Dm cross-domain, O open set, C continual. Robustness (Axis E): R any test, S social-network processing, A adversarial. \cmark{} = tested; -- = not tested; ? = unclear from the report. Str = strength of generalization evidence (VH very high, H high, M moderate, L low). QA = total appraisal score over QA1--QA8 (maximum 8); QA1, QA3 and QA8 were scored with AI assistance and validated by the authors (see the Quality Appraisal subsection of the main paper).

{\scriptsize
\setlength{\tabcolsep}{3pt}
\begin{longtable}{l c p{5.6cm} l l l c c c c c c c c c c c}
\caption{Coding of the included studies}\label{tab:coding-all}\\
\toprule
Study & Year & Short title & Task & Evid. & Gen. & U & F & Ds & Dm & O & C & R & S & A & Str & QA \\
\midrule
\endfirsthead
\toprule
Study & Year & Short title & Task & Evid. & Gen. & U & F & Ds & Dm & O & C & R & S & A & Str & QA \\
\midrule
\endhead
\bottomrule
\endfoot
\cite{R0001} R0001 & 2026 & A Multimodal Vision–language Model for Shrimp & Det & SEM & D & \cmark & -- & -- & -- & -- & -- & -- & -- & -- & H & 4.5 \\
\cite{R0002} R0002 & 2026 & PhyTrace & Det & GEO & Multi & \cmark & \cmark & \cmark & \cmark & -- & -- & \cmark & \cmark & -- & H & 7.0 \\
\cite{R0003} R0003 & 2024 & Robust GAN-Based CNN Model as Generative AI Ap & Det & SPA & G & -- & -- & -- & -- & -- & -- & -- & -- & -- & L & 3.0 \\
\cite{R0004} R0004 & 2026 & A Multi-Branch Transformer Based Cross-Attenti & Det & HYB & G+D & \cmark & \cmark & -- & -- & -- & -- & \cmark & -- & -- & M & 6.5 \\
\cite{R0005} R0005 & 2026 & Synthetic Image Detection Based on Multi-Level & Det & HYB & G+D & \cmark & \cmark & -- & -- & -- & -- & -- & -- & -- & M & 6.0 \\
\cite{R0006} R0006 & 2026 & SCA-Det & Det & FRQ & Multi & \cmark & \cmark & \cmark & \cmark & -- & -- & \cmark & -- & -- & VH & 6.5 \\
\cite{R0008} R0008 & 2026 & MS-DCL & Det & HYB & G+D & \cmark & \cmark & \cmark & \cmark & -- & -- & \cmark & -- & -- & VH & 6.5 \\
\cite{R0009} R0009 & 2026 & FMACNN & Det & SPA & D & -- & -- & -- & -- & -- & -- & \cmark & -- & -- & L & 5.5 \\
\cite{R0010} R0010 & 2024 & IPD-Net & Det & SPA & Multi & \cmark & \cmark & \cmark & -- & -- & -- & \cmark & -- & -- & VH & 6.0 \\
\cite{R0012} R0012 & 2026 & Penny-Wise and Pound-Foolish in AI-Generated I & Det & SEM & Multi & \cmark & \cmark & \cmark & \cmark & -- & -- & -- & -- & -- & H & 6.0 \\
\cite{R0013} R0013 & 2026 & Vulnerabilities in AI-Generated Image Detectio & Rob & FRQ & G+D & \cmark & \cmark & -- & -- & -- & -- & \cmark & -- & \cmark & M & 6.5 \\
\cite{R0014} R0014 & 2026 & A robust synthetic face detector in OSN contex & Det & SPA & G+D & \cmark & -- & -- & -- & -- & -- & \cmark & \cmark & -- & M & 6.5 \\
\cite{R0016} R0016 & 2025 & Transfer Learning-Driven Hybrid Architecture f & Det & HYB & G & \cmark & -- & \cmark & -- & -- & -- & -- & -- & -- & M & 5.0 \\
\cite{R0017} R0017 & 2024 & Generalized and robust model for GAN-generated & Det & NOI & G & \cmark & -- & -- & \cmark & -- & -- & \cmark & -- & -- & H & 6.0 \\
\cite{R0018} R0018 & 2026 & DCNet & Det & SPA & G+D & \cmark & \cmark & \cmark & \cmark & -- & -- & \cmark & -- & -- & H & 6.5 \\
\cite{R0019} R0019 & 2022 & An Eyes-Based Siamese Neural Network for the D & Det & SPA & G & \cmark & -- & -- & -- & -- & -- & \cmark & -- & -- & H & 6.5 \\
\cite{R0020} R0020 & 2026 & SF-CFNet & Det & HYB & G+D & ? & -- & \cmark & \cmark & -- & -- & \cmark & -- & -- & M & 6.0 \\
\cite{R0023} R0023 & 2026 & Cross-Model Deepfake Detection Through Contour & Det+Attr & FRQ & G+D & \cmark & \cmark & \cmark & -- & -- & -- & \cmark & -- & -- & M & 6.0 \\
\cite{R0027} R0027 & 2025 & Ensemble-Based Biometric Verification & Det & HYB & Multi & -- & -- & -- & -- & -- & -- & -- & -- & -- & L & 5.0 \\
\cite{R0029} R0029 & 2025 & Detection of AI generated images using combine & Rob & HYB & G+D & \cmark & \cmark & \cmark & \cmark & \cmark & -- & \cmark & -- & \cmark & M & 6.0 \\
\cite{R0031} R0031 & 2026 & Revisiting DIRE & Det & REC & Multi & \cmark & \cmark & \cmark & \cmark & -- & -- & \cmark & -- & -- & VH & 6.5 \\
\cite{R0032} R0032 & 2026 & SDA-Det & Det & SPA & Multi & \cmark & \cmark & \cmark & \cmark & -- & -- & \cmark & -- & -- & VH & 6.5 \\
\cite{R0034} R0034 & 2026 & BM-DDFN & Attr & REC & Multi & \cmark & -- & \cmark & \cmark & \cmark & -- & \cmark & \cmark & -- & M & 7.0 \\
\cite{R0035} R0035 & 2024 & A siamese-based verification system for open-s & Attr & SRC & Multi & \cmark & \cmark & -- & \cmark & \cmark & -- & -- & -- & -- & H & 6.5 \\
\cite{R0037} R0037 & 2024 & Detecting images generated by diffusers & Det & SEM & D & \cmark & -- & -- & -- & -- & -- & -- & -- & -- & M & 6.0 \\
\cite{R0039} R0039 & 2024 & Boosting Deep Feature Fusion-Based Detection M & Det & SPA & G & -- & -- & -- & -- & -- & -- & -- & -- & -- & L & 4.5 \\
\cite{R0040} R0040 & 2026 & Transformer Based on Multi-Domain Feature Fusi & Det & HYB & Multi & ? & ? & -- & -- & -- & -- & \cmark & -- & -- & L & 3.0 \\
\cite{R0041} R0041 & 2026 & VeriTrace & Bench & SPA & D & \cmark & -- & \cmark & \cmark & -- & -- & -- & -- & -- & M & 5.5 \\
\cite{R0045} R0045 & 2026 & Boon or Bane & Attr & SPA & D & -- & -- & -- & -- & -- & -- & \cmark & -- & -- & L & 5.5 \\
\cite{R0046} R0046 & 2025 & DiffCoR & Det & REC & Multi & \cmark & \cmark & -- & \cmark & -- & -- & -- & -- & -- & H & 6.0 \\
\cite{R0047} R0047 & 2024 & Artifact feature purification for cross-domain & Det & FRQ & G+D & \cmark & \cmark & -- & \cmark & -- & -- & -- & -- & -- & H & 5.5 \\
\cite{R0051} R0051 & 2021 & VIPPrint & Bench & SPA & G & -- & -- & -- & -- & -- & -- & \cmark & -- & -- & L & 5.5 \\
\cite{R0052} R0052 & 2025 & BOSC & Attr & SRC & Multi & \cmark & \cmark & \cmark & -- & \cmark & -- & \cmark & -- & -- & H & 7.0 \\
\cite{R0053} R0053 & 2025 & Advanced Detection of AI-Generated Images Thro & Det & SPA & G & -- & -- & -- & -- & -- & -- & -- & -- & -- & L & 4.5 \\
\cite{R0055} R0055 & 2026 & SIDWA & Det & HYB & Multi & \cmark & \cmark & ? & -- & -- & -- & \cmark & -- & -- & H & 6.5 \\
\cite{R0056} R0056 & 2026 & Background-Noise-Driven Detection of Diffusion & Det & NOI & Multi & \cmark & -- & \cmark & -- & -- & -- & -- & -- & -- & H & 5.5 \\
\cite{R0058} R0058 & 2026 & CLIP-ADA & Det & SEM & Multi & \cmark & \cmark & \cmark & \cmark & -- & -- & \cmark & -- & -- & VH & 6.0 \\
\cite{R0062} R0062 & 2024 & Enhancing Deepfake Detection With Diversified & Det & HYB & G & \cmark & -- & \cmark & -- & -- & -- & \cmark & -- & -- & M & 7.0 \\
\cite{R0063} R0063 & 2026 & A hybrid transfer learning approach for multi- & Det+Attr & HYB & G+D & -- & -- & -- & -- & -- & -- & -- & -- & -- & L & 5.5 \\
\cite{R0064} R0064 & 2025 & Multi-modal texture fusion network for detecti & Det & HYB & Multi & \cmark & \cmark & \cmark & -- & -- & -- & \cmark & -- & \cmark & H & 6.5 \\
\cite{R0065} R0065 & 2026 & AI-generated image detection using Machine Lea & Bench & HYB & Multi & \cmark & \cmark & -- & -- & -- & -- & -- & -- & -- & M & 6.5 \\
\cite{R0066} R0066 & 2025 & Vision Transformer-Based Framework for AI-Gene & Det & SPA & Multi & -- & -- & -- & -- & -- & -- & \cmark & -- & -- & L & 4.5 \\
\cite{R0067} R0067 & 2026 & An enhanced deep learning framework for DeepFa & Det & SPA & Multi & \cmark & ? & \cmark & ? & -- & -- & -- & -- & -- & M & 5.0 \\
\cite{R0069} R0069 & 2025 & Sanitizing Diffusion-Generated Images via Fing & Rob & FRQ & D & -- & -- & -- & -- & -- & -- & \cmark & -- & \cmark & L & 6.0 \\
\cite{R0070} R0070 & 2024 & Synthetic Face Discrimination via Learned Imag & Det & REC & G+D & \cmark & \cmark & \cmark & -- & -- & -- & \cmark & -- & -- & H & 7.0 \\
\cite{R0071} R0071 & 2024 & Less is more & Det & FRQ & G & -- & -- & \cmark & -- & -- & -- & \cmark & -- & -- & M & 7.0 \\
\cite{R0072} R0072 & 2026 & Geometric-semantic dual-adaptation for general & Det & SEM & Multi & \cmark & \cmark & -- & \cmark & -- & -- & -- & -- & -- & VH & 6.0 \\
\cite{R0073} R0073 & 2025 & Building a Universal Detector of AI-Generated & Det & HYB & Multi & \cmark & \cmark & \cmark & \cmark & -- & -- & \cmark & -- & \cmark & H & 7.0 \\
\cite{R0076} R0076 & 2025 & Benchmark dataset and framework for continual & Det+Attr & SEM & Multi & -- & -- & -- & -- & -- & \cmark & -- & -- & -- & L & 6.0 \\
\cite{R0077} R0077 & 2025 & MFF-Net & Det & HYB & Multi & \cmark & \cmark & \cmark & \cmark & -- & -- & \cmark & -- & -- & VH & 7.0 \\
\cite{R0078} R0078 & 2022 & EXPOSING AI GENERATED DEEPFAKE IMAGES USING SI & Det & SPA & G & \cmark & -- & \cmark & -- & -- & -- & -- & -- & -- & L & 6.5 \\
\cite{R0080} R0080 & 2026 & Structure searchable network model for deepfak & Det & HYB & Multi & \cmark & \cmark & \cmark & -- & -- & -- & \cmark & -- & -- & M & 7.0 \\
\cite{R0081} R0081 & 2024 & MMGANGuard & Det & SPA & G & -- & -- & -- & -- & -- & -- & -- & -- & -- & L & 4.5 \\
\cite{R0082} R0082 & 2026 & Bias-Free? An Empirical Study on Ethnicity, Ge & Bench & SPA & Multi & \cmark & \cmark & \cmark & -- & -- & -- & -- & -- & -- & H & 6.0 \\
\cite{R0083} R0083 & 2026 & DTBF & Det & HYB & G+D & \cmark & \cmark & \cmark & -- & -- & -- & -- & -- & -- & H & 5.0 \\
\cite{R0085} R0085 & 2024 & Performance Comparison and Visualization of AI & Bench & SPA & Multi & \cmark & \cmark & \cmark & \cmark & -- & -- & \cmark & -- & -- & M & 6.5 \\
\cite{R0086} R0086 & 2026 & Linguistic profiling of deepfakes & Bench & SEM & Multi & -- & -- & -- & -- & -- & -- & \cmark & -- & -- & L & 5.0 \\
\cite{R0089} R0089 & 2025 & Fusing Frequency and Spatial Transformers for & Det & HYB & G & \cmark & -- & -- & -- & -- & -- & -- & -- & -- & M & 6.0 \\
\cite{R0090} R0090 & 2026 & LiteUpdate & Det & HYB & G+D & \cmark & \cmark & -- & -- & -- & \cmark & \cmark & -- & -- & M & 6.5 \\
\cite{R0091} R0091 & 2026 & Generalizable and Adaptive Continual Learning & Det & SEM & Multi & \cmark & \cmark & \cmark & \cmark & -- & \cmark & \cmark & -- & -- & H & 5.5 \\
\cite{R0095} R0095 & 2025 & DeepGuard & Other & SRC & D & -- & -- & -- & -- & \cmark & -- & -- & -- & -- & L & 6.0 \\
\cite{R0096} R0096 & 2026 & Harnessing attention for cropping and fusion i & Det & HYB & Multi & \cmark & \cmark & \cmark & -- & \cmark & -- & \cmark & -- & -- & H & 5.5 \\
\cite{R0097} R0097 & 2026 & NSFF & Det & HYB & Multi & \cmark & \cmark & \cmark & -- & \cmark & -- & \cmark & -- & -- & M & 7.0 \\
\cite{R0098} R0098 & 2026 & A Transfer-Learning Approach for Detection of & Other & HYB & Multi & -- & -- & -- & -- & -- & -- & -- & -- & -- & L & 5.5 \\
\cite{R0099} R0099 & 2025 & Photo Authenticity Detection Using Machine Lea & Det & HYB & Multi & \cmark & -- & \cmark & \cmark & \cmark & -- & -- & -- & -- & M & 4.0 \\
\cite{R0101} R0101 & 2023 & Black-box attack against GAN-generated image d & Other & SPA & G & \cmark & -- & \cmark & -- & \cmark & -- & \cmark & -- & \cmark & M & 7.0 \\
\cite{R0104} R0104 & 2026 & ADRD & Det & REC & Multi & \cmark & \cmark & \cmark & \cmark & \cmark & -- & \cmark & -- & -- & M & 7.0 \\
\cite{R0105} R0105 & 2024 & Hierarchical Progressive Image Forgery Detecti & Det+Loc & HYB & Multi & \cmark & -- & \cmark & ? & -- & -- & -- & -- & -- & L & 5.0 \\
\cite{R0106} R0106 & 2025 & SCADET & Det & FRQ & Multi & \cmark & \cmark & -- & ? & -- & -- & \cmark & -- & -- & H & 7.0 \\
\cite{R0108} R0108 & 2023 & AI vs. AI & Det & SPA & Multi & \cmark & -- & \cmark & \cmark & -- & -- & -- & -- & -- & M & 6.5 \\
\cite{R0109} R0109 & 2026 & Lightweight AI-Generated image detection based & Det & SPA & Multi & \cmark & \cmark & \cmark & \cmark & -- & -- & \cmark & -- & -- & H & 7.0 \\
\cite{R0113} R0113 & 2026 & Dual-path AI-generated image detection & Det & HYB & Multi & \cmark & \cmark & \cmark & \cmark & -- & -- & \cmark & \cmark & -- & H & 7.5 \\
\cite{R0114} R0114 & 2024 & Reliable Out-of-Distribution Recognition of Sy & Det & FRQ & Multi & \cmark & \cmark & \cmark & \cmark & -- & -- & \cmark & \cmark & -- & H & 7.5 \\
\cite{R0115} R0115 & 2024 & Unsupervised Generative Fake Image Detector & Det & SPA & G+D & \cmark & \cmark & -- & -- & -- & -- & \cmark & -- & \cmark & M & 7.0 \\
\cite{R0118} R0118 & 2026 & Beyond Causal Models & Det & SPA & Multi & \cmark & \cmark & \cmark & -- & -- & -- & -- & -- & -- & H & 5.5 \\
\cite{R0120} R0120 & 2025 & Forensic Analysis of Manipulated Images and Vi & Bench & HYB & Multi & \cmark & \cmark & \cmark & -- & -- & -- & -- & -- & -- & M & 6.0 \\
\cite{R0121} R0121 & 2026 & Multi-Domain Perception Transformer for Genera & Det & HYB & Multi & \cmark & \cmark & \cmark & \cmark & -- & -- & -- & -- & -- & H & 6.5 \\
\cite{R0122} R0122 & 2025 & Detecting AI-Generated Images Using a Hybrid R & Det & HYB & D & -- & -- & -- & -- & -- & -- & -- & -- & -- & L & 5.5 \\
\cite{R0126} R0126 & 2024 & Synthbuster & Det & FRQ & D & \cmark & -- & \cmark & \cmark & \cmark & -- & \cmark & -- & -- & M & 7.0 \\
\cite{R0128} R0128 & 2025 & GCS-Net & Det & HYB & Multi & \cmark & \cmark & \cmark & \cmark & -- & -- & -- & -- & -- & VH & 6.5 \\
\cite{R0131} R0131 & 2025 & Enhancing deepfake detection with Adaptive-DCG & Det & HYB & G & -- & -- & -- & -- & -- & -- & -- & -- & -- & L & 5.0 \\
\cite{R0133} R0133 & 2026 & EXACT & Det+Attr & HYB & Multi & \cmark & \cmark & \cmark & \cmark & -- & -- & \cmark & -- & -- & H & 7.0 \\
\cite{R0135} R0135 & 2026 & Diffusion-Generated Image Detection via Target & Det & REC & G+D & \cmark & \cmark & \cmark & \cmark & -- & -- & -- & -- & -- & M & 6.5 \\
\cite{R0136} R0136 & 2026 & Integrating semantic structure and texture cue & Det & HYB & D & -- & -- & -- & -- & -- & -- & \cmark & -- & -- & L & 5.0 \\
\cite{R0144} R0144 & 2026 & Generalizable Synthetic Image Detection via La & Det & SEM & Multi & \cmark & \cmark & \cmark & \cmark & -- & -- & \cmark & \cmark & \cmark & VH & 7.0 \\
\cite{R0147} R0147 & 2026 & Spatial flatness-curvature mask driven general & Det & SPA & Multi & \cmark & \cmark & \cmark & \cmark & -- & -- & \cmark & -- & -- & H & 6.0 \\
\cite{R0148} R0148 & 2026 & DIFC-Net & Det & HYB & Multi & \cmark & \cmark & \cmark & \cmark & -- & -- & -- & -- & -- & H & 6.5 \\
\cite{R0152} R0152 & 2025 & Novel Deepfake Image Detection with PV-ISM & Det & SPA & Multi & -- & -- & \cmark & \cmark & -- & -- & -- & -- & -- & M & 5.5 \\
\cite{R0154} R0154 & 2025 & RADAR & Det & HYB & D & -- & -- & -- & \cmark & -- & -- & -- & -- & -- & M & 6.0 \\
\cite{R0155} R0155 & 2026 & Color Space Analysis for AI-Generated Image De & Det & SPA & G+D & \cmark & \cmark & \cmark & -- & -- & -- & \cmark & -- & -- & L & 6.5 \\
\cite{R0228} R0228 & 2026 & Proto-LeakNet & Attr & SPA & Multi & \cmark & \cmark & \cmark & \cmark & \cmark & -- & \cmark & \cmark & -- & H & 7.5 \\
\cite{R0242} R0242 & 2026 & The PRISM benchmark & Other & HYB & Multi & \cmark & \cmark & \cmark & \cmark & -- & -- & \cmark & -- & -- & H & 6.5 \\
\cite{R0259} R0259 & 2025 & How Much is the Source Mismatch an Important P & Other & SPA & Multi & \cmark & \cmark & \cmark & \cmark & \cmark & -- & \cmark & -- & -- & M & 6.5 \\
\cite{R0275} R0275 & 2022 & FD2Foremer & Other & HYB & Multi & \cmark & \cmark & \cmark & -- & -- & -- & \cmark & -- & -- & M & 7.0 \\
\cite{R0302} R0302 & 2026 & Toward Generalizable Forgery Detection and Rea & Det & HYB & Multi & \cmark & \cmark & \cmark & \cmark & -- & -- & \cmark & -- & -- & VH & 6.5 \\
\cite{R0309} R0309 & 2024 & Detection of real-time deep fakes and face for & Other & SPA & G & \cmark & -- & \cmark & -- & -- & -- & -- & -- & -- & M & 6.0 \\
\cite{R0321} R0321 & 2025 & DeepFeatureX-SN & Det+Attr & SPA & Multi & \cmark & -- & \cmark & -- & -- & -- & \cmark & -- & -- & M & 7.0 \\
\cite{R0325} R0325 & 2025 & DeepfakeCLIP & Other & SEM & G+D & \cmark & \cmark & \cmark & \cmark & \cmark & -- & \cmark & -- & -- & H & 6.0 \\
\cite{R0345} R0345 & 2023 & Parameter agnostic stacked wavelet transformer & Other & FRQ & Multi & \cmark & \cmark & \cmark & \cmark & \cmark & -- & \cmark & -- & \cmark & H & 7.0 \\

\end{longtable}
}

\section{Reported Performance Values}
\label{app:results}
Every performance value used in the Reported Performance Under Shift subsection of the main paper, as transcribed into workbook sheet~32 (\NresValues values from \NresStudies studies). Setting: in-domain, ood (out-of-distribution), open-set, robustness. Value in percent (values reported on a 0--1 scale are multiplied by 100). Stat: \emph{mean} if the study reports an average over several targets, otherwise \emph{value}. The workbook stores, for each row, the verbatim fragment of the extraction record from which the value was read. The authors validated every value.

{\scriptsize
\setlength{\tabcolsep}{4pt}
\begin{longtable}{l l l p{5.2cm} l r l}
\caption{Transcribed performance values}\label{tab:results-all}\\
\toprule
Study & Setting & Shift & Target & Metric & Value & Stat \\
\midrule
\endfirsthead
\toprule
Study & Setting & Shift & Target & Metric & Value & Stat \\
\midrule
\endhead
\bottomrule
\endfoot
% AUTO-GENERATED from sheet 32_Reported_Results (each value is quoted verbatim in the workbook)
\cite{R0001} & in-domain & none & same-pool split & AUC & 100.00 & value \\
\cite{R0001} & ood & unseen-generator & withheld Midjourney v6 & AUC & 98.90 & value \\
\cite{R0002} & ood & multiple & ImageNet benchmark & AUC & 81.98 & value \\
\cite{R0003} & in-domain & none & same-source split & ACC & 96.35 & value \\
\cite{R0004} & in-domain & none & GenImage SD1.4 & ACC & 99.15 & value \\
\cite{R0004} & robustness & jpeg/blur mix & GenImage SD1.4 & ACC & 98.36 & mean \\
\cite{R0005} & in-domain & none & GenImage Midjourney & ACC & 98.40 & value \\
\cite{R0005} & in-domain & none & GenImage Midjourney & AP & 99.50 & value \\
\cite{R0005} & ood & unseen-generator & GenImage (7 unseen) & ACC & 79.50 & mean \\
\cite{R0005} & ood & unseen-generator & GenImage (7 unseen) & AP & 86.60 & mean \\
\cite{R0006} & in-domain & none & ProGAN & ACC & 99.90 & value \\
\cite{R0006} & in-domain & none & ProGAN & AP & 100.00 & value \\
\cite{R0006} & ood & multiple & 45 generated-image sets & ACC & 93.00 & mean \\
\cite{R0006} & ood & multiple & 45 generated-image sets & AP & 98.40 & mean \\
\cite{R0006} & robustness & jpeg & JPEG Q30 & ACC & 88.30 & value \\
\cite{R0008} & in-domain & none & GenImage SD1.4 & ACC & 99.90 & value \\
\cite{R0008} & in-domain & none & GenImage SD1.4 & AP & 100.00 & value \\
\cite{R0008} & ood & unseen-generator & GenImage & ACC & 84.50 & mean \\
\cite{R0008} & ood & unseen-generator & GenImage & AP & 97.30 & mean \\
\cite{R0009} & in-domain & none & RealAIGI/CIFAKE/FF++ splits & ACC & 99.12 & value \\
\cite{R0010} & in-domain & none & ProGAN & ACC & 99.94 & value \\
\cite{R0010} & ood & cross-dataset & GenImage & ACC & 82.71 & mean \\
\cite{R0012} & ood & multiple & 10 benchmarks & ACC & 72.00 & mean \\
\cite{R0012} & ood & multiple & 10 benchmarks & AP & 86.14 & mean \\
\cite{R0014} & robustness & jpeg & JPEG Q50 & ACC & 98.29 & value \\
\cite{R0016} & in-domain & none & Kaggle source & ACC & 97.93 & value \\
\cite{R0016} & in-domain & none & Kaggle source & AUC & 98.80 & value \\
\cite{R0016} & ood & cross-dataset & FaceForensics++ & ACC & 94.18 & value \\
\cite{R0016} & ood & cross-dataset & Celeb-DF v2 & ACC & 91.87 & value \\
\cite{R0016} & ood & cross-dataset & DFDC & ACC & 89.45 & value \\
\cite{R0016} & ood & cross-dataset & FaceForensics++ & AUC & 97.00 & value \\
\cite{R0016} & ood & cross-dataset & Celeb-DF v2 & AUC & 96.10 & value \\
\cite{R0016} & ood & cross-dataset & DFDC & AUC & 94.90 & value \\
\cite{R0017} & in-domain & none & ProGAN (1-class) & ACC & 98.00 & value \\
\cite{R0017} & in-domain & none & ProGAN (1-class) & AP & 96.37 & value \\
\cite{R0017} & ood & unseen-generator & cross-GAN & ACC & 80.83 & mean \\
\cite{R0017} & ood & unseen-generator & cross-GAN & AP & 75.15 & mean \\
\cite{R0018} & ood & multiple & 22 generators & ACC & 90.78 & mean \\
\cite{R0019} & in-domain & none & StyleGAN2 & AUC & 100.00 & value \\
\cite{R0019} & ood & unseen-generator & ProGAN & AUC & 99.70 & value \\
\cite{R0019} & ood & unseen-generator & StyleGAN3 & AUC & 99.60 & value \\
\cite{R0020} & in-domain & none & SMTD & F1 & 98.96 & value \\
\cite{R0023} & ood & cross-dataset & leave-one-dataset-out & ACC & 75.91 & mean \\
\cite{R0023} & ood & cross-dataset & leave-one-dataset-out & AUC & 90.85 & mean \\
\cite{R0027} & in-domain & none & pooled six sources & ACC & 79.20 & value \\
\cite{R0029} & in-domain & none & GenImage SD1.5 & ACC & 99.80 & value \\
\cite{R0029} & ood & unseen-generator & Midjourney & ACC & 66.02 & value \\
\cite{R0029} & ood & unseen-generator & GLIDE & ACC & 73.61 & value \\
\cite{R0029} & ood & unseen-generator & VQDM & ACC & 57.17 & value \\
\cite{R0029} & ood & unseen-generator & BigGAN & ACC & 55.97 & value \\
\cite{R0029} & ood & unseen-generator & StyleGAN3 & ACC & 45.34 & value \\
\cite{R0031} & ood & multiple & 30 generators & ACC & 96.50 & mean \\
\cite{R0031} & ood & multiple & 30 generators & AP & 98.80 & mean \\
\cite{R0031} & ood & cross-dataset & Chameleon & ACC & 62.10 & value \\
\cite{R0032} & ood & multiple & 45 sets & ACC & 94.30 & mean \\
\cite{R0032} & ood & multiple & 45 sets & AP & 99.20 & mean \\
\cite{R0032} & robustness & jpeg & JPEG Q30 & ACC & 86.90 & value \\
\cite{R0034} & in-domain & none & AIGC-Trace closed set & ACC & 98.80 & value \\
\cite{R0034} & ood & cross-dataset & ForenSynths & ACC & 97.30 & value \\
\cite{R0034} & open-set & open-set & near-OOD rejection & AUC & 58.30 & value \\
\cite{R0035} & in-domain & none & closed set & AUC & 100.00 & value \\
\cite{R0035} & open-set & open-set & fully open (config 1) & AUC & 92.00 & value \\
\cite{R0035} & open-set & open-set & fully open (config 2) & AUC & 81.00 & value \\
\cite{R0035} & open-set & open-set & fully open (config 3) & AUC & 85.00 & value \\
\cite{R0039} & in-domain & none & 80/20 split & ACC & 91.32 & value \\
\cite{R0040} & in-domain & none & DFDC & ACC & 98.10 & value \\
\cite{R0046} & ood & unseen-generator & GenImage (SD1.4 protocol) & ACC & 77.50 & mean \\
\cite{R0047} & ood & unseen-generator & GenImage & ACC & 80.40 & mean \\
\cite{R0047} & ood & unseen-generator & DiffusionForensics & ACC & 98.70 & mean \\
\cite{R0051} & in-domain & none & digital validation & ACC & 100.00 & value \\
\cite{R0051} & robustness & print-and-scan & printed-and-scanned & ACC & 50.00 & value \\
\cite{R0052} & open-set & open-set & unknown architectures & AUC & 93.58 & mean \\
\cite{R0053} & in-domain & none & validation split & ACC & 98.20 & value \\
\cite{R0055} & in-domain & none & SIDset & ACC & 95.74 & value \\
\cite{R0055} & ood & unseen-generator & 9 external generators & ACC & 84.60 & mean \\
\cite{R0056} & in-domain & none & ADM & ACC & 98.00 & value \\
\cite{R0056} & in-domain & none & ADM & AP & 99.90 & value \\
\cite{R0056} & ood & unseen-generator & earlier generators & ACC & 89.90 & mean \\
\cite{R0056} & ood & unseen-generator & earlier generators & AP & 97.70 & mean \\
\cite{R0056} & ood & unseen-generator & recent generators & ACC & 93.60 & mean \\
\cite{R0056} & ood & unseen-generator & recent generators & AP & 99.30 & mean \\
\cite{R0058} & ood & cross-dataset & Synthbuster & ACC & 85.76 & value \\
\cite{R0058} & ood & multiple & diverse unseen sets & ACC & 75.82 & value \\
\cite{R0062} & ood & cross-dataset & Celeb-DF & AUC & 94.25 & value \\
\cite{R0062} & ood & cross-dataset & DFDC & AUC & 78.36 & value \\
\cite{R0063} & in-domain & none & skin-image dataset & ACC & 94.30 & value \\
\cite{R0064} & in-domain & none & ProGAN & ACC & 100.00 & value \\
\cite{R0064} & ood & unseen-generator & BigGAN & ACC & 95.20 & value \\
\cite{R0064} & ood & unseen-generator & CycleGAN & ACC & 98.50 & value \\
\cite{R0064} & ood & cross-family & GenImage Midjourney & ACC & 88.20 & value \\
\cite{R0064} & ood & cross-family & GenImage SD1.4 & ACC & 86.50 & value \\
\cite{R0064} & ood & cross-family & GenImage SD1.5 & ACC & 92.10 & value \\
\cite{R0064} & ood & cross-family & GenImage Wukong & ACC & 86.30 & value \\
\cite{R0064} & robustness & jpeg & BigGAN, JPEG Q50 & ACC & 89.10 & value \\
\cite{R0064} & robustness & adversarial & BigGAN, adversarial & ACC & 83.40 & value \\
\cite{R0066} & in-domain & none & validation split & ACC & 96.25 & value \\
\cite{R0067} & in-domain & none & internal test & ACC & 97.95 & value \\
\cite{R0071} & in-domain & none & FFHQ/StyleGAN2 & ACC & 99.20 & value \\
\cite{R0071} & ood & cross-dataset & CelebA + StyleGAN2 & ACC & 99.11 & mean \\
\cite{R0072} & in-domain & none & AIGIBench same-source & ACC & 99.70 & mean \\
\cite{R0072} & ood & unseen-generator & GenImage (SD1.4 protocol) & ACC & 96.30 & mean \\
\cite{R0073} & ood & unseen-generator & UniversalFakeDetect-style benchmark & ACC & 95.58 & mean \\
\cite{R0077} & in-domain & none & ProGAN & AP & 100.00 & value \\
\cite{R0077} & ood & multiple & 18 datasets & AP & 96.13 & mean \\
\cite{R0078} & in-domain & none & FFHQ/StyleGAN & ACC & 94.80 & value \\
\cite{R0078} & ood & unseen-generator & PGAN & ACC & 54.75 & value \\
\cite{R0080} & in-domain & none & seen sources & ACC & 95.04 & value \\
\cite{R0080} & in-domain & none & seen sources & AUC & 98.99 & value \\
\cite{R0080} & ood & cross-dataset & ten unseen subsets & ACC & 65.90 & mean \\
\cite{R0080} & ood & cross-dataset & ten unseen subsets & AUC & 77.46 & mean \\
\cite{R0080} & robustness & jpeg & JPEG Q50 & ACC & 93.92 & value \\
\cite{R0083} & ood & multiple & 16 generators & ACC & 96.13 & mean \\
\cite{R0083} & ood & multiple & 16 generators & AP & 99.15 & mean \\
\cite{R0086} & in-domain & none & DFLIP-3K & ACC & 92.30 & value \\
\cite{R0086} & robustness & jpeg & JPEG Q50 & ACC & 89.00 & value \\
\cite{R0089} & in-domain & none & validation & ACC & 87.32 & value \\
\cite{R0089} & ood & unseen-generator & StyleGAN3 & ACC & 93.00 & value \\
\cite{R0095} & in-domain & none & DeepGuardDB & ACC & 99.87 & value \\
\cite{R0096} & in-domain & none & ProGAN & ACC & 100.00 & value \\
\cite{R0096} & ood & multiple & 16 generators & ACC & 94.40 & mean \\
\cite{R0096} & robustness & jpeg & JPEG & ACC & 83.70 & value \\
\cite{R0097} & ood & multiple & 14 categories & ACC & 80.82 & mean \\
\cite{R0097} & robustness & jpeg & JPEG & ACC & 79.55 & value \\
\cite{R0098} & in-domain & none & custom test set & ACC & 98.50 & value \\
\cite{R0099} & in-domain & none & curated collection & ACC & 94.80 & value \\
\cite{R0104} & in-domain & none & SD v1.5 self-test & ACC & 71.80 & value \\
\cite{R0104} & in-domain & none & SD v1.5 self-test & AP & 78.50 & value \\
\cite{R0104} & ood & unseen-generator & leave-one-model-out & ACC & 61.20 & mean \\
\cite{R0104} & ood & unseen-generator & leave-one-model-out & AP & 65.10 & mean \\
\cite{R0105} & in-domain & none & AITfake & ACC & 99.56 & mean \\
\cite{R0105} & ood & cross-dataset & CoCoGLIDE & ACC & 54.19 & value \\
\cite{R0105} & ood & cross-dataset & HIFI-IFDL & ACC & 82.97 & value \\
\cite{R0105} & ood & cross-dataset & GenImage & ACC & 80.28 & value \\
\cite{R0105} & ood & cross-dataset & CASIA & ACC & 77.16 & value \\
\cite{R0106} & in-domain & none & AI-ArtBench & AUC & 96.20 & value \\
\cite{R0108} & in-domain & none & RSI test & ACC & 100.00 & value \\
\cite{R0108} & ood & cross-dataset & ADE20K & ACC & 88.90 & value \\
\cite{R0108} & ood & cross-dataset & Sketchy & ACC & 94.30 & value \\
\cite{R0108} & ood & cross-dataset & CUB-200-2011 & ACC & 97.80 & value \\
\cite{R0109} & in-domain & none & ProGAN & ACC & 99.95 & value \\
\cite{R0109} & in-domain & none & ProGAN & AP & 100.00 & value \\
\cite{R0109} & ood & multiple & AIGCDetectBenchmark (16) & ACC & 90.97 & mean \\
\cite{R0109} & ood & multiple & AIGCDetectBenchmark (16) & AP & 96.62 & mean \\
\cite{R0113} & in-domain & none & UFD Benchmark 1 & ACC & 99.50 & value \\
\cite{R0113} & in-domain & none & UFD Benchmark 1 & AP & 99.90 & value \\
\cite{R0113} & ood & unseen-generator & UFD Benchmark 1 OOD & ACC & 93.00 & mean \\
\cite{R0113} & ood & unseen-generator & UFD Benchmark 1 OOD & AP & 97.30 & mean \\
\cite{R0113} & ood & cross-dataset & Chameleon & ACC & 67.83 & value \\
\cite{R0113} & ood & cross-dataset & Synthbuster & ACC & 78.18 & mean \\
\cite{R0113} & robustness & jpeg & JPEG QF50 & ACC & 82.79 & value \\
\cite{R0114} & in-domain & none & in-distribution test & AUC & 99.30 & value \\
\cite{R0114} & ood & multiple & uncompressed OOD & ACC & 70.40 & mean \\
\cite{R0118} & ood & unseen-generator & GAN dataset & ACC & 98.80 & mean \\
\cite{R0118} & ood & cross-family & diffusion dataset & ACC & 96.30 & mean \\
\cite{R0121} & in-domain & none & Celeb-DF & ACC & 96.30 & value \\
\cite{R0122} & in-domain & none & CIFAKE & ACC & 96.12 & value \\
\cite{R0126} & in-domain & none & generic model & MCC & 87.20 & value \\
\cite{R0126} & ood & unseen-generator & generalization training & MCC & 52.70 & value \\
\cite{R0128} & in-domain & none & ProGAN & ACC & 99.88 & value \\
\cite{R0128} & in-domain & none & ProGAN & AP & 100.00 & value \\
\cite{R0128} & ood & multiple & 19 unseen sets & ACC & 84.57 & mean \\
\cite{R0128} & ood & multiple & 19 unseen sets & AP & 92.46 & mean \\
\cite{R0131} & in-domain & none & 20\% split & ACC & 95.00 & value \\
\cite{R0133} & in-domain & none & in-domain test & ACC & 99.49 & value \\
\cite{R0133} & ood & cross-dataset & 11 cross-dataset domains & ACC & 84.88 & mean \\
\cite{R0133} & robustness & jpeg & JPEG Q30 & ACC & 81.42 & value \\
\cite{R0135} & in-domain & none & ADM (LSUN) & ACC & 82.70 & value \\
\cite{R0135} & ood & unseen-generator & 4 diffusion generators & ACC & 80.40 & mean \\
\cite{R0135} & ood & cross-family & StyleGAN & ACC & 66.00 & value \\
\cite{R0136} & in-domain & none & custom split & ACC & 98.50 & value \\
\cite{R0144} & ood & unseen-generator & Tgan & AP & 97.40 & mean \\
\cite{R0147} & in-domain & none & ProGAN & ACC & 98.70 & value \\
\cite{R0147} & in-domain & none & ProGAN & AP & 99.90 & value \\
\cite{R0147} & ood & multiple & 38 unseen generators & ACC & 87.20 & mean \\
\cite{R0147} & ood & multiple & 38 unseen generators & AP & 94.60 & mean \\
\cite{R0148} & in-domain & none & SD v1.5 & AUC & 99.02 & value \\
\cite{R0148} & ood & unseen-generator & 10 unseen generators & AUC & 90.29 & mean \\
\cite{R0152} & in-domain & none & CIFAKE & ACC & 96.60 & value \\
\cite{R0155} & in-domain & none & CIFAKE & ACC & 98.91 & value \\
\cite{R0228} & in-domain & none & WILD closed set & AUC & 98.13 & value \\
\cite{R0228} & open-set & open-set & WILD open set & AUC & 100.00 & value \\
\cite{R0242} & in-domain & none & PRISM full generation & ACC & 96.60 & mean \\
\cite{R0242} & ood & cross-dataset & GenImage & ACC & 90.30 & mean \\
\cite{R0242} & ood & cross-dataset & DRCT-2M & ACC & 97.40 & mean \\
\cite{R0275} & ood & cross-dataset & Celeb-DF & AUC & 79.20 & value \\
\cite{R0275} & ood & cross-dataset & DFDC & AUC & 73.60 & value \\
\cite{R0302} & ood & unseen-generator & MMFR (10 generators) & ACC & 94.06 & mean \\
\cite{R0302} & ood & cross-dataset & LOKI & ACC & 86.73 & value \\
\cite{R0321} & in-domain & none & held-out split (binary) & ACC & 97.70 & value \\
\cite{R0321} & ood & multiple & generalization sets & ACC & 67.40 & mean \\
\cite{R0321} & robustness & jpeg & JPEG QF50 & ACC & 69.85 & value \\
\cite{R0325} & in-domain & none & ProGAN & ACC & 99.50 & value \\
\cite{R0325} & in-domain & none & ProGAN & AP & 100.00 & value \\
\cite{R0325} & ood & unseen-generator & cross-GAN & ACC & 92.90 & mean \\
\cite{R0325} & ood & unseen-generator & cross-GAN & AP & 98.50 & mean \\
\cite{R0325} & ood & cross-family & cross-diffusion & ACC & 89.50 & mean \\
\cite{R0325} & ood & cross-family & cross-diffusion & AP & 97.70 & mean \\

\end{longtable}
}

\section{Full-Text Exclusions and Reports Not Retrieved}
\label{app:excluded}
Table~S3 lists the \NexclFT reports excluded at full-text assessment, with the exclusion code defined in the Review Methodology section of the main paper. Table~S4 lists the \Nnotretrieved reports that were sought but could not be retrieved through institutional access.

{\small
\begin{longtable}{l p{12.5cm} l}
\caption{Reports excluded at full text}\label{tab:ftex}\\
\toprule
Record & Title & Code \\
\midrule
\endfirsthead
\toprule
Record & Title & Code \\
\midrule
\endhead
\bottomrule
\endfoot
% AUTO-GENERATED
R0028 & Training humans for synthetic face image detection & FT-EC5 \\
R0061 & Design of an integrated model combining vision transformer, multi-modal fusion, & FT-EC9 \\
R0123 & Generative Adversarial Ensemble Learning for Face Forensics & FT-EC4 \\
R0130 & ADVANCED DEEPFAKE DETECTION LEVERAGING SWIN TRANSFORMER TECHNOLOGY & FT-EC4 \\
R0151 & Hybrid Deep Learning for Image Authenticity: Distinguishing Between Real and AI- & FT-EC3 \\
R0217 & Deepfake Video Forensics: A Hybrid Deep Learning Model for Detecting Deepfake Fa & FT-EC14 \\
R0329 & An efficient deepfake video detection using robust deep learning & FT-EC14 \\

\end{longtable}

\begin{longtable}{l p{14.2cm} l}
\caption{Reports sought but not retrieved}\label{tab:notret}\\
\toprule
Record & Title & Year \\
\midrule
\endfirsthead
\toprule
Record & Title & Year \\
\midrule
\endhead
\bottomrule
\endfoot
% AUTO-GENERATED
R0007 & AI-generated image detection algorithm based on classical-quantum hybrid neural network & 2026 \\
R0015 & ImageAuthNet: toward an explainable and generalizable framework to detect fake images - use cas & 2026 \\
R0022 & DADP: a discrepancy-aware dual-path network for AI-generated image detection & 2026 \\
R0026 & Semantic Distribution and Authenticity Discrepancy Alignment for AI-Generated Image Detection & 2026 \\
R0043 & Image quality assessment based fake face detection & 2023 \\
R0054 & Bi-LORA: A Vision-Language Approach for Synthetic Image Detection & 2025 \\
R0059 & Research on adversarial identification methods for AI-generated image software Craiyon V3 & 2025 \\
R0060 & Self-Supervised AI-Generated Image Detection: A Camera Metadata Perspective & 2026 \\
R0075 & ArtGate: Injecting Fake Artifact Features into CLIP for AI-Generated Image Detection & 2026 \\
R0079 & Cross-generator image authenticity verification: multi-modal feature fusion and distribution di & 2026 \\
R0084 & SPRITZ-PS: validation of synthetic face images using a large dataset of printed documents & 2024 \\
R0087 & Toward Generalizable AI-Generated Image Detection with a New Realistic Dataset: Performance Eva & 2026 \\
R0088 & Source Attribution of AI-Generated Images in Data Scarcity Conditions via Training-Free Resynth & 2026 \\
R0100 & High-resolution network-based multi-feature fusion for generalized forgery detection & 2025 \\
R0107 & AI-Generated Image Detection Method Integrating Light -Shadow Sensitive Features and Kolmogorov & 2025 \\
R0112 & MSAFNet: multi-scale self-adaptive feature fusion network for AI-generated image detection & 2025 \\
R0116 & Synthetic image detection method integrating neural ordinary differential equations and spatial & 2026 \\
R0117 & Transferable Dual-Domain Feature Importance Attack Against AI-Generated Image Detector & 2026 \\
R0119 & ERASE: Bypassing Collaborative Detection of AI Counterfeit via Comprehensive Artifacts Eliminat & 2026 \\
R0125 & CoDA: Color Distribution Probing for Efficient and Generalizable AI-Generated Image Detection & 2026 \\
R0127 & Hybrid AI approaches for detecting deepfake faces & 2025 \\
R0129 & Quality-Guided Forgery Adapter for Generalizable AIGC Image Detection & 2026 \\
R0132 & Multi-Domain Feature Fusion: Leveraging Color, Noise, and Spatial Features for AI-generated Ima & 2026 \\
R0134 & Adversarial Mixture of Experts Post-Training for Robust AI-Generated Image Detection & 2026 \\
R0138 & Dual Frequency Branch Framework With Reconstructed Sliding Windows Attention for AI-Generated I & 2026 \\
R0139 & FALCON-Net: Feature Aggregation of Local Patterns for AI-Generated Image Detection & 2026 \\
R0140 & Dual-scale model collaborative reasoning with multi-feature fusion for robust AI-generated imag & 2026 \\
R0141 & Towards Extensible Detection of AI-Generated Images via Content-Agnostic Adapter-Based Category & 2025 \\
R0142 & Optimized Frequency Collaborative Strategy Drives AI Image Detection & 2025 \\
R0143 & Enhancing deep fake detection: comparative analysis of advanced deep learning architectures & 2026 \\
R0145 & FAMSeC: A Few-Shot-Sample-Based General AI-Generated Image Detection Method & 2025 \\
R0149 & Wavelet-packets for deepfake image analysis and detection & 2022 \\
R0150 & Learning a Semantic Similarity Orthogonal Space for Model-Level AI-Generated Image Source Attri & 2026 \\
R0153 & Incremental Learning for AI-Generated Image Detection & 2026 \\
R0222 & FD-GAN: Generalizable and Robust Forgery Detection via Generative Adversarial Networks & 2024 \\
R0252 & Beyond Familiar Domains: A Study of the Generalization Ability of Machine-Generated Image Detec & 2025 \\
R0265 & RealDTT: Towards A Comprehensive Real-World Dataset for Tampered Text Detection & 2025 \\
R0301 & Exploring Bi-Level Inconsistency via Blended Images for Generalizable Face Forgery Detection & 2024 \\
R0339 & Enhancing Deepfake Detection: Leveraging StyleGAN3 for Robust AI-Generated Forgery Identificati & 2026 \\

\end{longtable}
}
